% ===== PRÉAMBULE CANONIQUE — à coller VERBATIM en tête de CHAQUE .tex =====
\documentclass[11pt,a4paper]{article}
\usepackage[utf8]{inputenc}\usepackage[T1]{fontenc}\usepackage{lmodern}
\usepackage[french]{babel}
\usepackage{amsmath,amssymb,mathtools}\usepackage{icomma}
\usepackage[version=4]{mhchem}          % \ce{...} : équations & formules
\usepackage{siunitx}\sisetup{locale=FR} % \qty{}{}, \si{}, \ang{}
\usepackage{chemfig}                    % \chemfig{...} : structures développées
\usepackage{xcolor}\usepackage{colortbl}
\usepackage{tikz}\usepackage{pgfplots}\pgfplotsset{compat=1.18}
\usepackage[most]{tcolorbox}\usepackage{enumitem}\usepackage{array,booktabs}
\usepackage[a4paper,margin=2.2cm]{geometry}
\usetikzlibrary{arrows.meta,calc,angles,quotes,positioning,patterns,backgrounds,shapes.geometric}
\definecolor{primary}{RGB}{222,49,99}\definecolor{primarydark}{RGB}{150,22,55}
\definecolor{defcol}{RGB}{0,114,178}\definecolor{methcol}{RGB}{52,92,148}
\definecolor{excol}{RGB}{0,135,90}\definecolor{attncol}{RGB}{200,40,55}
\tcbset{boxbase/.style={enhanced,breakable,boxrule=0.9pt,arc=2.5pt,
  left=3mm,right=3mm,top=2.3mm,bottom=2.3mm,fonttitle=\bfseries,coltitle=white}}
% --- boîtes TUTORIEL (noms EXACTS, ne pas renommer) ---
\newtcolorbox{cle}[1][]{boxbase,colback=defcol!6,colframe=defcol,
  title={\textsc{À retenir}\ifx\relax#1\relax\relax\else\textmd{ : \itshape #1}\fi}}
\newtcolorbox{methode}[1][]{boxbase,colback=methcol!7,colframe=methcol,sharp corners,
  title={\textsc{Méthode}\ifx\relax#1\relax\relax\else\textmd{ --- \itshape #1}\fi}}
\newtcolorbox{exemplebox}[1][]{boxbase,colback=excol!5,colframe=excol,
  title={\textsc{Exemple}\ifx\relax#1\relax\relax\else\textmd{ : \itshape #1}\fi}}
\newtcolorbox{piege}[1][]{boxbase,colback=attncol!6,colframe=attncol,
  title={\textsc{Piège}\ifx\relax#1\relax\relax\else\textmd{ : \itshape #1}\fi}}
\newtcolorbox{remarque}[1][]{boxbase,colback=black!4,colframe=black!45,
  title={\textsc{Remarque}\ifx\relax#1\relax\relax\else\textmd{ : \itshape #1}\fi}}
% --- boîtes EXERCICE / CORRIGÉ (noms EXACTS, auto-comptées) ---
\newtcolorbox[auto counter]{exo}[1][]{boxbase,colback=primary!4,colframe=primary,
  title={\textsc{Exercice}~\thetcbcounter\ifx\relax#1\relax\relax\else\textmd{ --- \itshape #1}\fi}}
\newtcolorbox[auto counter]{corrige}[1][]{boxbase,colback=excol!4,colframe=excol,
  title={\textsc{Corrigé}~\thetcbcounter\ifx\relax#1\relax\relax\else\textmd{ --- \itshape #1}\fi}}
\newcommand{\R}{\mathbb{R}}\newcommand{\N}{\mathbb{N}}\newcommand{\Z}{\mathbb{Z}}\newcommand{\ds}{\displaystyle}
% (un \usepackage{fancyhdr}+en-tête est possible ; il est ignoré côté web)
% =========================================================================
\begin{document}

\begin{exo}[une même pression dans quatre unités]
Une pression vaut $p=\qty{2}{atm}$. Exprime-la successivement dans les unités demandées (on prendra
$\qty{1}{atm}=\qty{101325}{\pascal}=\qty{760}{mmHg}=\qty{1.013}{bar}$).
\begin{enumerate}[label=\alph*)]
  \item en \si{\pascal} ;
  \item en \si{\kilo\pascal} ;
  \item en \si{bar} ;
  \item en \si{mmHg}.
\end{enumerate}
\end{exo}

\begin{exo}[température et rapport de températures]
On chauffe un gaz de $\theta_1=\qty{27}{\degreeCelsius}$ à $\theta_2=\qty{87}{\degreeCelsius}$.
\begin{enumerate}[label=\alph*)]
  \item Donne $T_1$ en kelvins.
  \item Donne $T_2$ en kelvins.
  \item Calcule le rapport $T_2/T_1$ (avec les kelvins).
  \item Calcule le « rapport » $\theta_2/\theta_1$ avec les degrés Celsius ; qu'en conclus-tu ?
\end{enumerate}
\end{exo}

\begin{exo}[du volume d'une pièce au volume de gaz]
Une salle a pour dimensions $\qty{4}{\meter}\times\qty{5}{\meter}\times\qty{2.5}{\meter}$.
\begin{enumerate}[label=\alph*)]
  \item Calcule son volume en \si{\meter\cubed}.
  \item Convertis ce volume en \si{\litre}.
  \item Convertis $\qty{22.4}{\meter\cubed}$ en \si{\litre}.
  \item Convertis $\qty{0.25}{\meter\cubed}$ en \si{\litre}.
\end{enumerate}
\end{exo}

\begin{exo}[lecture d'un manomètre]
Un manomètre indique une pression $p=\qty{1520}{mmHg}$.
\begin{enumerate}[label=\alph*)]
  \item Exprime $p$ en \si{atm}.
  \item Exprime $p$ en \si{bar}.
  \item Exprime $p$ en \si{\kilo\pascal}.
  \item Combien de fois cette pression dépasse-t-elle la pression atmosphérique ($\qty{1}{atm}$) ?
\end{enumerate}
\end{exo}

\end{document}
